% Focus: Move technical computations supporting Sections 5--6 here; keep the main text clean and theorem-driven.

\subsection{Contact computations and bracket identities}\label{subsec:appB-contact}
% Focus: Provide full computations of $\alpha\wedge d\alpha$, horizontal frames, and commutator identities.

We collect the computations used in Section~\ref{subsec:contact-form}--\ref{subsec:contact-horizontal}.

\subsubsection*{Contact condition}
With $\omega=\mu\,du+\lambda a\,dv$ and $\alpha=da-\omega$, we have
\[
d\omega=\lambda\,da\wedge dv,
\qquad
d\alpha=-\lambda\,da\wedge dv.
\]
Therefore
\[
\alpha\wedge d\alpha
=(da-\mu\,du-\lambda a\,dv)\wedge(-\lambda\,da\wedge dv)
=\mu\lambda\,du\wedge da\wedge dv.
\]
In particular, $\alpha$ is a contact form if and only if $\mu\lambda\neq 0$.

\subsubsection*{Horizontal frame and bracket}
The horizontal lifts $D_u=\partial_u+\mu\partial_a$ and $D_v=\partial_v+\lambda a\,\partial_a$ satisfy $\alpha(D_u)=\alpha(D_v)=0$. Their bracket is
\[
[D_u,D_v]=[\partial_u+\mu\partial_a,\partial_v+\lambda a\partial_a]=\mu\lambda\,\partial_a.
\]

\subsubsection*{Darboux coordinates}
Rescale $\tilde u=\mu u$ and $\tilde v=\lambda v$, and set $x=\tilde v$, $y=a$, $z=a-\tilde u$. Then $dz=da-\mu\,du$ and $dx=\lambda\,dv$, so
\[
\alpha=da-\mu\,du-\lambda a\,dv=dz-y\,dx,
\]
which is the Darboux normal form.

\subsection{Details of the \texorpdfstring{$\delta$}{delta}-calculus}\label{subsec:appB-delta}
% Focus: Derive formulas for $\delta$, its algebraic properties, and its relation to $d$ on horizontal forms.

For a smooth scalar field $F(u,v,a)$,
\[
\delta F=(D_uF)\,du+(D_vF)\,dv.
\]
Using $da=\alpha+\omega$ with $\omega=\mu\,du+\lambda a\,dv$, we obtain
\[
dF-(\partial_aF)\alpha
=(\partial_uF+\mu\partial_aF)\,du+(\partial_vF+\lambda a\,\partial_aF)\,dv
=(D_uF)\,du+(D_vF)\,dv
=\delta F.
\]
Moreover,
\[
\delta^2F=([D_u,D_v]F)\,du\wedge dv=\mu\lambda(\partial_aF)\,du\wedge dv.
\]

\subsection{Model holomorphicity computations}\label{subsec:appB-holomorphic}
% Focus: Collect explicit calculations for the first-order system and the associated twisted second-order equations.

We record the algebra underlying Proposition~\ref{prop:analytic-factorization} and Theorem~\ref{thm:analytic-twisted-harmonicity}.

For complex-valued $F$, define
\[
\bar\partial_{\mathrm{AEG}}=\frac12(D_u+iD_v),
\qquad
\partial_{\mathrm{AEG}}=\frac12(D_u-iD_v).
\]
Then
\[
4\,\partial_{\mathrm{AEG}}\bar\partial_{\mathrm{AEG}}F
=(D_u-iD_v)(D_u+iD_v)F
=(D_u^2+D_v^2)F+i(D_uD_v-D_vD_u)F.
\]
Using $[D_u,D_v]=\mu\lambda\,\partial_a$, we obtain
\[
4\,\partial_{\mathrm{AEG}}\bar\partial_{\mathrm{AEG}}F
=\bigl(\Delta_H+i\mu\lambda\,\partial_a\bigr)F,
\qquad
\Delta_H=D_u^2+D_v^2.
\]
Similarly,
\[
4\,\bar\partial_{\mathrm{AEG}}\partial_{\mathrm{AEG}}F
=\bigl(\Delta_H-i\mu\lambda\,\partial_a\bigr)F.
\]
If $F$ is horizontally holomorphic, i.e.\ $\bar\partial_{\mathrm{AEG}}F=0$, then applying $\partial_{\mathrm{AEG}}$ yields
\[
\bigl(\Delta_H+i\mu\lambda\,\partial_a\bigr)F=0,
\]
which is the twisted harmonicity equation.
